\documentclass[12pt, a4paper]{article}

\usepackage{fontspec}
\usepackage{polyglossia}
\usepackage{graphicx}
\setmainlanguage{russian}
\setotherlanguage{english}
\setmainfont{DejaVu Serif}
\setsansfont{DejaVu Sans}
\setmonofont{DejaVu Sans Mono}

\usepackage[left=25mm, right=20mm, top=20mm, bottom=20mm]{geometry}
\usepackage{booktabs}
\usepackage{tabularx}
\usepackage{array}
\usepackage{hyperref}
\usepackage{parskip}
\usepackage{enumitem}
\usepackage{titlesec}

\hypersetup{
    colorlinks=true,
    linkcolor=black,
    urlcolor=blue
}

\titleformat{\section}{\large\bfseries}{}{0em}{}[\titlerule]
\titleformat{\subsection}{\normalsize\bfseries}{}{0em}{}

\setlength{\parindent}{0pt}
\setlength{\parskip}{6pt}

\begin{document}

% =================== ТИТУЛЬНЫЙ ЛИСТ ===================
\thispagestyle{empty}

\begin{center}

{\small
Московский ордена Трудового Красного Знамени\\
Физико-технический институт\\
(Национальный исследовательский университет)\\
Факультет аэрофизики и космических исследований
}

\vspace{1cm}

\vfill

{\Huge\bfseries ASTRA}

\vspace{0.5cm}

{\Large Автоматизированный лазерный звёздный трекер}

\vspace{0.5cm}

{\normalsize Проектное предложение}

\vfill

{\small Долгопрудный, 2026}

\end{center}

\newpage

% =================== ОСНОВНОЕ СОДЕРЖАНИЕ ===================
\setcounter{page}{1}

\section*{Название проекта}

\textbf{ASTRA} (Autonomous Stellar Targeting Remote-Access Apparatus) --- автоматизированный лазерный звёздный трекер для наведения на небесные объекты.

\section*{Ментор}

Провоторов Павел Владимирович

\section*{Команда}

\begin{tabularx}{\textwidth}{|X|c|c|c|}
\hline
\textbf{Участник} & \textbf{Группа} & \textbf{Telegram} & \textbf{Телефон} \\
\hline
Градовцев Святослав Алексеевич & Б03-503 & @svt\_gradov & +7~999~248-86-97 \\
\hline
Коломоец Александр Анатольевич & Б03-503 & @AlexanderKolomoets & +7~961~213-93-21 \\
\hline
Минаева Полина Игоревна & Б03-502 & @Polina\_Minaewa & +7~910~855-06-48 \\
\hline
\end{tabularx}

\section*{Цель проекта}

Создать автоматизированное устройство для наведения лазерного указателя (5~мВт, 532~нм) на любую звезду, планету или иной небесный объект по выбору пользователя, обеспечивающее полный охват небосвода по азимуту и возвышению (360° по каждой оси без мёртвых зон), автономную работу не менее 3~часов и точность наведения 0{,}5--2° после калибровки по звёздам, с управлением через Telegram-бот по Wi-Fi.

\section*{Описание функционала}

Компактный звёздный трекер на двухосной поворотной платформе. Устройство получает данные о местоположении пользователя через GPS-модуль, определяет ориентацию относительно горизонта при помощи гироскопа и акселерометра, и на основе астрономических расчётов наводит лазер на любой выбранный небесный объект.

\subsection*{Основные функции}

\begin{itemize}[leftmargin=*, nosep]
    \item Автоматическое определение координат пользователя (GPS NEO-6M)
    \item Определение угла наклона и ориентации устройства (IMU MPU-6050)
    \item Расчёт положения небесных тел в реальном времени
    \item Наведение лазера на выбранный объект по двум осям
    \item Управление через Telegram-бот по Wi-Fi
    \item Калибровка устройства перед использованием (датчики Холла, процедура star alignment)
\end{itemize}

\subsection*{Дополнительный функционал}

\begin{itemize}[leftmargin=*, nosep]
    \item Режим автоматического сопровождения объекта
    \item Встроенный каталог {\textasciitilde}200 объектов во flash ESP32 + Yale Bright Star Catalog (9110~звёзд) в SPIFFS
    \item Ручной ввод координат как fallback при отсутствии GPS-сигнала
    \item RTC DS3231 для точного хранения времени при потере GPS
\end{itemize}

\subsection*{Целевые характеристики}

\begin{tabularx}{\textwidth}{|X|X|}
\hline
\textbf{Параметр} & \textbf{Значение} \\
\hline
Диапазон азимута & 360° непрерывно \\
\hline
Диапазон возвышения & 360° непрерывно \\
\hline
Точность наведения & 0{,}5--2°  \\
\hline
Угловое разрешение & 0{,}0028°/шаг \\
\hline
Лазер & Зелёный 532~нм, 5~мВт, класс 3R \\
\hline
Управление & Telegram-бот по Wi-Fi (ESP32) \\
\hline
Питание & АКБ 3S Li-Ion 11{,}1~В / 3000~мАч \\
\hline
Автономность & 3 часа \\
\hline
Масса & {\textasciitilde}1{,}1--1{,}3~кг \\
\hline
Габариты & {\textasciitilde}150$\times$150$\times$190~мм (цилиндр) + основание Ø220~мм \\
\hline
Корпус & 3D-печать PLA/PETG \\
\hline
\end{tabularx}

\section*{Конструкция}

\subsection*{Форм-фактор}

Цилиндрический корпус ({\textasciitilde}150$\times$150$\times$190~мм) на круглом основании (Ø220~мм), установленный в U-рамку на поворотной платформе.

\subsection*{Две оси вращения}

\textbf{Ось азимута --- нижнее основание:}
\begin{itemize}[leftmargin=*, nosep]
    \item Диапазон: 360° непрерывно
    \item Привод: NEMA17 + червячная пара 1:40
    \item Slip Ring \#1 в центре основания --- передаёт питание и сигналы при любом угле вращения
\end{itemize}

\textbf{Ось возвышения --- горизонтальная ось U-рамки:}
\begin{itemize}[leftmargin=*, nosep]
    \item Диапазон: 360° непрерывно
    \item Привод: NEMA17 + червячная пара 1:40
    \item Slip Ring \#2 в одной из опор --- передаёт питание и сигнал к лазеру
\end{itemize}

\section*{Задачи проекта}

\begin{enumerate}[leftmargin=*, nosep]
    \item Выбор элементной базы и архитектуры системы
    \item Разработка CAD-моделей корпуса и поворотной платформы (SolidWorks)
    \item Проектирование механизма наведения лазера (двухосная монтировка, червячные передачи)
    \item Создание электронной схемы устройства
    \item Разработка собственной печатной платы
    \item Интеграция GPS-модуля, гироскопа и управляющего контроллера
    \item Разработка прошивки микроконтроллера
    \item Реализация Telegram-бота
    \item Производство деталей методом 3D-печати и лазерной резки
    \item Сборка и тестирование устройства
    \item Калибровка системы и проверка точности наведения
\end{enumerate}

\section*{Анализ существующих аналогов}

\subsection*{1. Celestron NexStar SE (серия GoTo-телескопов)}

Линейка автоматизированных телескопов с computerized GoTo-монтировкой. 
Устройство автоматически наводится на выбранный объект из базы данных 
40\,000+ небесных тел и удерживает его в поле зрения. Выравнивание 
выполняется по технологии SkyAlign --- достаточно навестись на любые 
три ярких объекта.

\textbf{Преимущества:} полностью автоматическое наведение и сопровождение 
объектов, большая база данных, простое выравнивание, высокая точность 
червячных приводов.\\
\textbf{Недостатки:} очень высокая стоимость (от \$500), большие габариты 
и масса (требует штатива), необходимость телескопической оптики, 
отсутствие лазерного указателя для визуального наблюдения без окуляра.\\
\url{https://www.celestron.com/products/nexstar-5se-computerized-telescope}

\subsection*{2. DIY OpenAstroTracker}

Open-source проект монтировки для астрофотографии на Arduino. Устройство компенсирует вращение Земли, удерживая телескоп или камеру на выбранном объекте, однако не осуществляет автоматического наведения — ориентация выполняется вручную с выравниванием по Полярной звезде.

\textbf{Преимущества:} открытый исходный код, возможность самостоятельной сборки, низкая стоимость.\\
\textbf{Недостатки:} отсутствие автоматического наведения на объект, необходимость ручной полярной юстировки, наличие мёртвых зон по углу возвышения, не предназначен для лазерного указания.\\
\url{https://openastrotech.com}

\textbf{Отличие ASTRA от аналогов:} автономная работа с GPS и гироскопической системой, удобное управление через Telegram-бот, полный охват небосвода без мёртвых зон по обеим осям, компактный корпус с лазерным указателем вместо телескопической оптики.

\section*{Эскиз проекта}

Устройство состоит из цилиндрического корпуса с лазерным модулем, установленного в U-рамку на поворотной платформе. Нижнее основание содержит привод азимута и аккумулятор, U-рамка несёт привод возвышения и оба slip ring для непрерывного вращения по обеим осям.

\begin{figure}[h]
    \centering
    \includegraphics[width=0.9\textwidth]{Лазерный трекер.png}
    \caption{Эскиз устройства ASTRA}
\end{figure}

\section*{Элементная база}

\begin{tabularx}{\textwidth}{|X|X|c|}
\hline
\textbf{Компонент} & \textbf{Модель} & \textbf{Кол-во} \\
\hline
Микроконтроллер & ESP32 DevKit v1 & 1 \\
\hline
Драйвер шаговых двигателей & TMC2209 StepStick & 2 \\
\hline
Шаговый двигатель & NEMA17, 40~мм, 1{,}8°/шаг & 2 \\
\hline
GPS-модуль & NEO-6M с антенной & 1 \\
\hline
Гироскоп и акселерометр & MPU-6050 (IMU, I2C) & 1 \\
\hline
Часы реального времени & DS3231 (RTC, I2C) & 1 \\
\hline
Лазерный модуль & Зелёный 532~нм, 5~мВт, TTL & 1 \\
\hline
MOSFET-ключ лазера & IRLZ44N & 2 \\
\hline
Токосъёмное кольцо & Slip Ring 12~мм, 6~каналов & 2 \\
\hline
Аккумулятор & 18650 Li-Ion 3S, 3000~мАч + BMS & 1 \\
\hline
DC-DC преобразователь & MP1584, 11{,}1~В $\to$ 5~В & 2 \\
\hline
Подшипники & 608ZZ & 4 \\
\hline
Датчики нулевого положения & Датчик Холла A3144 & 2 \\
\hline
Печатная плата & Заказная PCB (KiCad $\to$ JLCPCB) & 1 \\
\hline
Корпусные элементы & PLA/PETG, 3D-печать, лазерная резка & --- \\
\hline
\end{tabularx}

\end{document}
